
\section{Définition}


\begin{frame}
\frametitle{Une définition (stricte) de la scalabilité}


\begin{defblock}{Définition}
Un \alert{système} est \alert{scalable}  si ses \alert{performances}
sont proportionnelles aux \alert{ressources} qui lui sont allouées.
\end{defblock}

où
\begin{redbullet}
 \item le \alert{système} est un système distribué de gestion de données;
  \item les \alert{ressources} sont les composants matériels (serveurs, connectique);
  \item la \alert{performance} est mesurée
  \begin{itemize}
    \item soit par la \alert{latence} pour les systèmes temps réel (temps d'accès à un document);
    \item soit par le \alert{débit} pour les systèmes analytiques (temps de parcours d'une collection).
  \end{itemize}
\end{redbullet}

On quantifie les grandeurs ``ressource'' et ``performance'' ; on montre que 
leur rapport est constant pour un volume de données fixé.

\end{frame}


\section{Exemple 1: comptons les documents}

\begin{frame}{Premier exemple: comptage des vidéos}

On veut compter les vidéos dans la collection. 

\vfill

\begin{redbullet}
  \item Avec un serveur: on lit la collection sur le disque, on compte: 100 MO/s
  \item Avec deux serveurs: on répartit la collection en deux parties égales, chaque
      serveur lit sur son disque: 200 MO/s
   \item avec $n$ serveurs: on répartit en $n$, on lit en parallèle: $n \times 100$ MO/s
\end{redbullet} 

À la fin, le ou les serveurs envoie(nt) leur compte à l'application qui cumule (coût 0 ou presque).

\vfill

Rapport performance/ressources = (débit global / nb serveurs) = 100 MO/s !
\end{frame}

\begin{frame}{La méthode}


\figSlideWithSize{scalability1}{10cm}

\vfill

\begin{remblock}{Important}
Le fait que le temps soit constant par serveur (courbe verte) est une condition nécessaire mais pas suffisante!
\end{remblock}

Si j'envoie \alert{toutes} les vidéos au client, le coût réseau (500 TOs!) devient prépondérant.
\end{frame}


\section{Exemple 2: cherchons les doublons}

\begin{frame}{Second exemple: on cherche des doublons}

Comment savoir si on a des vidéos en double (des \alert{contenus} identiques)?

\vfill

On procède en deux étapes

\begin{redbullet}

  \item parcours des vidéos, calcul d'une \alert{signature}; on produit la liste
  des paires $(i, s)$, $i$ id de la vidéo, $s$ signature;
  
  \item regroupement des paires par signature:  si deux signatures (ou plus) égales sont trouvées,
    c'est un doublon!
\end{redbullet}

Et on envoie au client la liste des doublons.
\end{frame}


\begin{frame}{Etape 1}

La première étape se déroule entièrement en parallèle sur les serveurs.

\vfill


\figSlideWithSize{scalability2}{10cm}


\vfill

Scalable, pour les mêmes raisons que l'exemple 1.
\end{frame}


\begin{frame}{Etape 2}

La  seconde étape implique des échanges réseaux: on alloue à chaque serveur un
sous-ensemble des signatures, et on lui transmet les paires de la liste pour ces signatures.

\vfill

\figSlideWithSize{scalability3}{10cm}


\vfill

Scalable? Peut-on augmenter linéairement le débit d'un réseau?

\end{frame}


\begin{frame}{Analyse}

On échange sur le réseau un paire $(i, s)$ pour chaque document.

\vfill

La taille est \alert{totalement négligeable} par rapport à celle du document lui-même.

\vfill

Conclusion: \alert{le coût des accès disques est largement prépondérant; la méthode est 
scalable (pas parfaitement...)}

\begin{remblock}{A vérifier}
Il faudrait \alert{quantifier} les coûts respectifs des accès disques et réseau (cf. la suite)
\end{remblock}

\end{frame}


\begin{frame}{Quelques conclusions générales}

Quelques conditions favorables (voire indispensables) 

\begin{redbullet}
  \item Données équitablement réparties.
  \item Traitement essentiellement \alert{local}.
  \item Limitation des transferts réseaux, entre les serveurs et vers le client.
  \item Et un temps de calcul global qui reste raisonnable pour les ressources qu'on peut/veut lui allouer...
\end{redbullet}

\begin{remblock}{Scalable mais pas acceptable?}
Le fait d'être scalable (selon notre définition) ne veut pas dire que le temps
de calcul est \alert{acceptable}. Si un traitement prend 10 ans, il prendra encore 5 ans en doublant les ressources...
\end{remblock}


\end{frame}

\end{document}


